%% kapa: exam=1 round=37 year=2026 period=1 form=A subject=economics subject_ko=경제학원론 seq=23 number=63
%% source: Q-Net 감정평가사 기출문제 (cst00309) artlSeq=5258924 file=raw/1008/5258924/2026년 제37회 감정평가사 1차 1교시.pdf
%% conversion: kordoc 4.24.51 -> tools/gen1.py (auto); review=proofread-sonnet; flags=none
\begin{question}{63}
소비이론에 관한 설명으로 옳지 않은 것은?
\begin{choices}
  \item 피셔(I. Fisher)의 두 기간 소비선택이론에 의하면 이자율의 변동은 정상재의 현재소비에 영향을 미치지 않는다.
  \item 예비적저축가설은 소득의 불확실성과 소비와의 관계를 설명한다.
  \item 케인즈의 절대소득가설에 의하면 일시적인 세율인하는 소비의 증가로 이어진다.
  \item 생애주기가설에 의하면 고령인구의 비율이 높아질수록 민간부문의 저축률이 하락한다.
  \item 항상소득가설에 의하면 현재의 소득이 항상소득과 일시소득의 합으로 결정된다.
\end{choices}
\end{question}
\answer{1}
